\section{财务交付总览}
\begin{exhibitbox}[表：2026Q1 业绩交付与全年折算]
\centering\tiny
\renewcommand{\arraystretch}{1.18}
\setlength{\tabcolsep}{2pt}
\begin{tabularx}{\exhibitboxwidth}{>{\bfseries\raggedright\arraybackslash}p{1.35cm} >{\raggedleft\arraybackslash}p{0.82cm} >{\raggedleft\arraybackslash}p{0.82cm} >{\raggedleft\arraybackslash}p{1.10cm} >{\raggedleft\arraybackslash}p{0.78cm} >{\raggedleft\arraybackslash}p{0.78cm} >{\raggedleft\arraybackslash}p{0.86cm} >{\raggedright\arraybackslash\sloppy\hspace{0pt}}X}
\toprule
\textbf{标的} & \textbf{Q1收入} & \textbf{Q1归母} & \textbf{收入/利润YoY} & \textbf{毛利率} & \textbf{OCF} & \textbf{26E归母} & \textbf{交付判断} \\
\midrule
中际旭创 300308 & 194.96 & 57.35 & \makecell[r]{+192.1\%\\+262.3\%} & 46.1\% & 33.68 & 249.33 & 强交付 \\
新易盛 300502 & 83.38 & 27.80 & \makecell[r]{+105.8\%\\+76.8\%} & 49.2\% & 6.84 & 120.88 & 稳定交付 \\
天孚通信 300394 & 13.30 & 4.92 & \makecell[r]{+40.8\%\\+45.8\%} & 56.6\% & 1.82 & 20.51 & 稳定交付 \\
光迅科技 002281 & 27.73 & 2.40 & \makecell[r]{+24.8\%\\+59.8\%} & 26.8\% & 1.96 & 10.91 & 稳定交付 \\
华工科技 000988 & 42.66 & 6.38 & \makecell[r]{+27.1\%\\+55.8\%} & 19.8\% & -11.27 & 25.54 & 稳定交付 \\
剑桥科技 603083 & 12.87 & 1.18 & \makecell[r]{+44.0\%\\+276.4\%} & 29.3\% & -2.19 & 5.92 & 强交付 \\
太辰光 300570 & 3.41 & 0.66 & \makecell[r]{-8.0\%\\-17.1\%} & 37.8\% & -0.93 & 2.63 & 承压 \\
源杰科技 688498 & 3.55 & 1.79 & \makecell[r]{+320.9\%\\+1153.1\%} & 77.8\% & 0.40 & 7.80 & 强交付 \\
亨通光电 600487 & 177.91 & 11.05 & \makecell[r]{+34.1\%\\+98.5\%} & 16.0\% & -1.94 & 46.06 & 稳定交付 \\
中天科技 600522 & 131.42 & 9.19 & \makecell[r]{+34.7\%\\+46.4\%} & 15.6\% & -19.49 & 38.29 & 稳定交付 \\
长飞光纤 601869 & 36.95 & 4.95 & \makecell[r]{+27.7\%\\+226.4\%} & 41.5\% & 6.18 & 21.53 & 强交付 \\
烽火通信 600498 & 45.10 & 0.38 & \makecell[r]{+11.4\%\\-30.4\%} & 26.4\% & -7.88 & 1.75 & 承压 \\
中兴通讯 000063 & 349.88 & 13.10 & \makecell[r]{+6.1\%\\-46.6\%} & 28.3\% & -19.79 & 59.57 & 承压 \\
仕佳光子 688313 & 5.77 & 1.16 & \makecell[r]{+32.2\%\\+24.7\%} & 34.1\% & -0.97 & 5.05 & 稳定交付 \\
长光华芯 688048 & 1.30 & 0.04 & \makecell[r]{+37.8\%\\+159.7\%} & 30.1\% & -1.31 & 0.22 & 强交付 \\
光库科技 300620 & 4.26 & 0.45 & \makecell[r]{+60.8\%\\+312.5\%} & 36.6\% & -0.34 & 1.95 & 强交付 \\
通鼎互联 002491 & 10.64 & 0.51 & \makecell[r]{+61.1\%\\-61.1\%} & 27.0\% & -0.18 & 2.23 & 承压 \\
永鼎股份 600105 & 12.46 & 1.59 & \makecell[r]{+41.9\%\\-45.2\%} & 26.2\% & 2.59 & 6.62 & 承压 \\
长芯博创 300548 & 6.71 & 1.30 & \makecell[r]{+24.5\%\\+45.0\%} & 49.2\% & 2.26 & 5.42 & 稳定交付 \\
德科立 688205 & 2.54 & 0.20 & \makecell[r]{+28.0\%\\+35.1\%} & 25.7\% & -0.64 & 0.89 & 稳定交付 \\
腾景科技 688195 & 1.71 & 0.14 & \makecell[r]{+51.2\%\\+10.7\%} & 37.2\% & -0.05 & 0.63 & 承压 \\
联特科技 301205 & 2.13 & 0.03 & \makecell[r]{-9.6\%\\-83.7\%} & 43.4\% & -1.08 & 0.15 & 承压 \\
兆龙互连 300913 & 5.08 & 0.49 & \makecell[r]{+13.5\%\\+49.0\%} & 21.1\% & 0.52 & 2.02 & 稳定交付 \\
意华股份 002897 & 12.92 & 0.40 & \makecell[r]{-5.3\%\\-37.7\%} & 20.0\% & 2.74 & 1.83 & 承压 \\
神宇股份 300563 & 2.04 & 0.13 & \makecell[r]{+16.7\%\\+11.4\%} & 18.2\% & 0.07 & 0.61 & 承压 \\
\bottomrule
\end{tabularx}
\end{exhibitbox}
\sourcenote{公司财务数据来自 astock.cli financials；2026E 归母净利使用公司层面 Q1 季节性假设折算，不等同于外部一致预期。}

\section{公司层面判断}
\textbf{光模块与器件层}：中际旭创和新易盛拥有最直接的 AI 模块利润弹性；长芯博创、德科立和联特科技补足相干传输、高速模块和高 beta 模块层；天孚通信、光库科技、太辰光和腾景科技更多体现精密器件、无源器件和封装价值；光迅科技、华工科技、剑桥科技具备平台和模块业务，但需要用 AI 数通收入占比与利润率证明业务纯度。

\textbf{上游光芯片/光源层}：源杰科技、长光华芯和仕佳光子分别映射高速激光芯片、激光器与 PLC/AWG/无源光芯片。它们的战略稀缺性更强，但估值容忍度必须绑定客户认证、良率、产能利用率和收入占比，而不是只绑定“国产替代”叙事。

\textbf{光纤光缆、线缆互连与网络设备层}：长飞光纤、亨通光电、中天科技、通鼎互联、永鼎股份提供预制棒、光纤、光缆、海缆和 ODN/通信线缆底座；兆龙互连、意华股份、神宇股份补足高速线缆、连接器和通信互连；中兴通讯、烽火通信提供运营商/云网络设备和应用层映射。这一层的估值更依赖运营商 capex、FTTH/DCI 项目、现金流、线缆价格和网络设备利润率，不能照搬模块龙头倍数。

\textbf{设备和材料观察池}：罗博特科及 MOCVD、外延、光刻/刻蚀、主动耦合、高速测试、InP/GaAs/薄膜铌酸锂/硅光晶圆等材料设备环节是完整产业链的一部分，但本报告不在缺少稳定利润分母时强行给目标价。它们进入观察池，后续若出现可验证盈利和纯度，再单独建模。

\section{覆盖标的逐项公司卡}
下面的公司卡将产业位置、2026Q1 财务交付、估值含义和下一季度验证点逐项展开。其目的不是重复估值表，而是把每一只股票为什么被给到该动作说清楚，避免用板块景气度覆盖个股差异。

\subsection{中际旭创（300308）：AI 数据中心高速光模块}
中际旭创 位于本报告定义的 \textbf{核心光模块} 层级，研究权重为 13\%。2026Q1 公司收入为 CNY194.96 亿元，归母净利润为 CNY57.35 亿元，毛利率 46.1\%，归母净利同比 +262.3\%。按照本报告季节性假设，2026E 归母净利为 CNY249.33 亿元，2027E EPS 为 CNY29.73。

\begin{exhibitbox}[单票卡：中际旭创]
\small
\begin{tabularx}{\exhibitboxwidth}{>{\bfseries\raggedright\arraybackslash}p{2.5cm} >{\raggedright\arraybackslash\sloppy\hspace{0pt}}X}
\toprule
产业位置 & AI 数据中心高速光模块 \\
当前价与市值 & CNY1253.89；市值约 CNY13881 亿元 \\
估值结论 & 综合目标 CNY1229.56，区间 CNY959--1781，空间 -1.9\%，动作 \stance{riskamber}{中性} \\
估值锚点 & 内在锚 CNY1337.79；市场锚 CNY1003.11；权重 内在65\% / 市场20\% / 券商15\% \\
核心催化 & 1.6T 放量、海外云厂商订单能见度提升、毛利率维持在 42\%以上。 \\
失效条件 & 800G/1.6T 订单放缓、客户集中度压力上升，或毛利率跌破 38\%。 \\
\bottomrule
\end{tabularx}
\end{exhibitbox}

投资含义上，800G/1.6T 规模和盈利兑现最强，但当前估值已经反映海外 AI 需求持续高景气。 市场锚的作用是承认当前成交和共识定价已经给出的情绪溢价，但它不能替代利润、毛利率和现金流。组合处理应当把它放在与当前动作匹配的仓位层级，而不是因为光通信行业景气就无差别买入。下一季度最关键的是利润、毛利率和现金流是否同步，若只有收入增长而现金流或毛利率恶化，则说明订单质量低于股价隐含预期。


\subsection{新易盛（300502）：高速光模块}
新易盛 位于本报告定义的 \textbf{核心光模块} 层级，研究权重为 12\%。2026Q1 公司收入为 CNY83.38 亿元，归母净利润为 CNY27.80 亿元，毛利率 49.2\%，归母净利同比 +76.8\%。按照本报告季节性假设，2026E 归母净利为 CNY120.88 亿元，2027E EPS 为 CNY15.83。

\begin{exhibitbox}[单票卡：新易盛]
\small
\begin{tabularx}{\exhibitboxwidth}{>{\bfseries\raggedright\arraybackslash}p{2.5cm} >{\raggedright\arraybackslash\sloppy\hspace{0pt}}X}
\toprule
产业位置 & 高速光模块 \\
当前价与市值 & CNY566.00；市值约 CNY5620 亿元 \\
估值结论 & 综合目标 CNY576.55，区间 CNY450--862，空间 +1.9\%，动作 \stance{riskamber}{中性} \\
估值锚点 & 内在锚 CNY633.04；市场锚 CNY452.80；权重 内在65\% / 市场20\% / 券商15\% \\
核心催化 & 海外客户订单延续、1.6T 产品认证、现金转化改善。 \\
失效条件 & 订单前置后回落、模块 ASP 下跌超过 20\%，或经营现金流明显落后利润。 \\
\bottomrule
\end{tabularx}
\end{exhibitbox}

投资含义上，高增长光模块弹性强，2025--2026 年交付较强，但估值对单一客户和 ASP 假设敏感。 市场锚的作用是承认当前成交和共识定价已经给出的情绪溢价，但它不能替代利润、毛利率和现金流。组合处理应当把它放在与当前动作匹配的仓位层级，而不是因为光通信行业景气就无差别买入。下一季度最关键的是利润、毛利率和现金流是否同步，若只有收入增长而现金流或毛利率恶化，则说明订单质量低于股价隐含预期。


\subsection{天孚通信（300394）：精密光器件与封装平台}
天孚通信 位于本报告定义的 \textbf{精密光器件} 层级，研究权重为 6\%。2026Q1 公司收入为 CNY13.30 亿元，归母净利润为 CNY4.92 亿元，毛利率 56.6\%，归母净利同比 +45.8\%。按照本报告季节性假设，2026E 归母净利为 CNY20.51 亿元，2027E EPS 为 CNY3.27。

\begin{exhibitbox}[单票卡：天孚通信]
\small
\begin{tabularx}{\exhibitboxwidth}{>{\bfseries\raggedright\arraybackslash}p{2.5cm} >{\raggedright\arraybackslash\sloppy\hspace{0pt}}X}
\toprule
产业位置 & 精密光器件与封装平台 \\
当前价与市值 & CNY318.00；市值约 CNY2472 亿元 \\
估值结论 & 综合目标 CNY254.40，区间 CNY198--293，空间 -20.0\%，动作 \stance{riskamber}{中性观察} \\
估值锚点 & 内在锚 CNY147.20；市场锚 CNY254.40；权重 内在57\% / 市场28\% / 券商15\% \\
核心催化 & 微光学器件和封装件在 1.6T/CPO 产品中的附加值提升。 \\
失效条件 & 器件降价、客户自供比例提高，或毛利率回落至 50\%以下。 \\
\bottomrule
\end{tabularx}
\end{exhibitbox}

投资含义上，高毛利精密器件平台价值明确，但增速相对模块龙头放缓。 市场锚的作用是承认当前成交和共识定价已经给出的情绪溢价，但它不能替代利润、毛利率和现金流。组合处理应当把它放在与当前动作匹配的仓位层级，而不是因为光通信行业景气就无差别买入。下一季度最关键的是利润、毛利率和现金流是否同步，若只有收入增长而现金流或毛利率恶化，则说明订单质量低于股价隐含预期。


\subsection{光迅科技（002281）：电信/数通光器件与光模块}
光迅科技 位于本报告定义的 \textbf{光器件/光模块} 层级，研究权重为 4\%。2026Q1 公司收入为 CNY27.73 亿元，归母净利润为 CNY2.40 亿元，毛利率 26.8\%，归母净利同比 +59.8\%。按照本报告季节性假设，2026E 归母净利为 CNY10.91 亿元，2027E EPS 为 CNY1.64。

\begin{exhibitbox}[单票卡：光迅科技]
\small
\begin{tabularx}{\exhibitboxwidth}{>{\bfseries\raggedright\arraybackslash}p{2.5cm} >{\raggedright\arraybackslash\sloppy\hspace{0pt}}X}
\toprule
产业位置 & 电信/数通光器件与光模块 \\
当前价与市值 & CNY240.34；市值约 CNY1922 亿元 \\
估值结论 & 综合目标 CNY187.47，区间 CNY146--216，空间 -22.0\%，动作 \stance{riskamber}{中性观察} \\
估值锚点 & 内在锚 CNY58.46；市场锚 CNY187.47；权重 内在47\% / 市场38\% / 券商15\% \\
核心催化 & 数通收入占比提高、硅光/CPO 产品验证推进。 \\
失效条件 & 电信周期走弱、数通结构未改善，或毛利率持续低于 28\%。 \\
\bottomrule
\end{tabularx}
\end{exhibitbox}

投资含义上，产品线和央企平台优势存在，但盈利能力显著低于 AI 光模块龙头。 市场锚的作用是承认当前成交和共识定价已经给出的情绪溢价，但它不能替代利润、毛利率和现金流。组合处理应当把它放在与当前动作匹配的仓位层级，而不是因为光通信行业景气就无差别买入。下一季度最关键的是利润、毛利率和现金流是否同步，若只有收入增长而现金流或毛利率恶化，则说明订单质量低于股价隐含预期。


\subsection{华工科技（000988）：光模块与激光设备混合平台}
华工科技 位于本报告定义的 \textbf{混合光通信平台} 层级，研究权重为 4\%。2026Q1 公司收入为 CNY42.66 亿元，归母净利润为 CNY6.38 亿元，毛利率 19.8\%，归母净利同比 +55.8\%。按照本报告季节性假设，2026E 归母净利为 CNY25.54 亿元，2027E EPS 为 CNY3.02。

\begin{exhibitbox}[单票卡：华工科技]
\small
\begin{tabularx}{\exhibitboxwidth}{>{\bfseries\raggedright\arraybackslash}p{2.5cm} >{\raggedright\arraybackslash\sloppy\hspace{0pt}}X}
\toprule
产业位置 & 光模块与激光设备混合平台 \\
当前价与市值 & CNY160.94；市值约 CNY1606 亿元 \\
估值结论 & 综合目标 CNY128.75，区间 CNY100--148，空间 -20.0\%，动作 \stance{riskamber}{中性观察} \\
估值锚点 & 内在锚 CNY92.79；市场锚 CNY128.75；权重 内在47\% / 市场38\% / 券商15\% \\
核心催化 & 高速光模块收入占比提升、工业激光盈利改善。 \\
失效条件 & 营运资本压力、混合业务折价扩大，或数通订单不放量。 \\
\bottomrule
\end{tabularx}
\end{exhibitbox}

投资含义上，混合业务提供盈利稳定性，但削弱纯 AI 光模块弹性。 市场锚的作用是承认当前成交和共识定价已经给出的情绪溢价，但它不能替代利润、毛利率和现金流。组合处理应当把它放在与当前动作匹配的仓位层级，而不是因为光通信行业景气就无差别买入。下一季度最关键的是利润、毛利率和现金流是否同步，若只有收入增长而现金流或毛利率恶化，则说明订单质量低于股价隐含预期。


\subsection{剑桥科技（603083）：光模块与通信设备高 beta 标的}
剑桥科技 位于本报告定义的 \textbf{高弹性主题} 层级，研究权重为 2\%。2026Q1 公司收入为 CNY12.87 亿元，归母净利润为 CNY1.18 亿元，毛利率 29.3\%，归母净利同比 +276.4\%。按照本报告季节性假设，2026E 归母净利为 CNY5.92 亿元，2027E EPS 为 CNY2.21。

\begin{exhibitbox}[单票卡：剑桥科技]
\small
\begin{tabularx}{\exhibitboxwidth}{>{\bfseries\raggedright\arraybackslash}p{2.5cm} >{\raggedright\arraybackslash\sloppy\hspace{0pt}}X}
\toprule
产业位置 & 光模块与通信设备高 beta 标的 \\
当前价与市值 & CNY245.92；市值约 CNY856 亿元 \\
估值结论 & 综合目标 CNY177.06，区间 CNY138--204，空间 -28.0\%，动作 \stance{riskamber}{中性观察} \\
估值锚点 & 内在锚 CNY105.68；市场锚 CNY177.06；权重 内在47\% / 市场38\% / 券商15\% \\
核心催化 & 800G/1.6T 订单持续转化、费用率被收入放量摊薄。 \\
失效条件 & Q2/Q3 利润不放量、现金流仍为负，或客户结构恶化。 \\
\bottomrule
\end{tabularx}
\end{exhibitbox}

投资含义上，反转弹性存在，但 Q1 利润分母仍偏小，难以支撑过高目标价。 市场锚的作用是承认当前成交和共识定价已经给出的情绪溢价，但它不能替代利润、毛利率和现金流。组合处理应当把它放在与当前动作匹配的仓位层级，而不是因为光通信行业景气就无差别买入。下一季度最关键的是利润、毛利率和现金流是否同步，若只有收入增长而现金流或毛利率恶化，则说明订单质量低于股价隐含预期。


\subsection{太辰光（300570）：光纤连接器与无源光器件}
太辰光 位于本报告定义的 \textbf{无源器件} 层级，研究权重为 2\%。2026Q1 公司收入为 CNY3.41 亿元，归母净利润为 CNY0.66 亿元，毛利率 37.8\%，归母净利同比 -17.1\%。按照本报告季节性假设，2026E 归母净利为 CNY2.63 亿元，2027E EPS 为 CNY1.33。

\begin{exhibitbox}[单票卡：太辰光]
\small
\begin{tabularx}{\exhibitboxwidth}{>{\bfseries\raggedright\arraybackslash}p{2.5cm} >{\raggedright\arraybackslash\sloppy\hspace{0pt}}X}
\toprule
产业位置 & 光纤连接器与无源光器件 \\
当前价与市值 & CNY234.10；市值约 CNY531 亿元 \\
估值结论 & 综合目标 CNY163.87，区间 CNY128--188，空间 -30.0\%，动作 \stance{riskamber}{中性观察} \\
估值锚点 & 内在锚 CNY26.60；市场锚 CNY163.87；权重 内在38\% / 市场47\% / 券商15\% \\
核心催化 & AI 数据中心连接器订单恢复、Q2 利润重新加速。 \\
失效条件 & 收入继续下滑、无源器件降价，或库存重建停滞。 \\
\bottomrule
\end{tabularx}
\end{exhibitbox}

投资含义上，无源器件敞口真实，但 2026Q1 收入/利润下滑，估值需要交付修复。 市场锚的作用是承认当前成交和共识定价已经给出的情绪溢价，但它不能替代利润、毛利率和现金流。组合处理应当把它放在与当前动作匹配的仓位层级，而不是因为光通信行业景气就无差别买入。下一季度最关键的是利润、毛利率和现金流是否同步，若只有收入增长而现金流或毛利率恶化，则说明订单质量低于股价隐含预期。


\subsection{源杰科技（688498）：激光芯片与高速光源}
源杰科技 位于本报告定义的 \textbf{战略上游} 层级，研究权重为 4\%。2026Q1 公司收入为 CNY3.55 亿元，归母净利润为 CNY1.79 亿元，毛利率 77.8\%，归母净利同比 +1153.1\%。按照本报告季节性假设，2026E 归母净利为 CNY7.80 亿元，2027E EPS 为 CNY12.54。

\begin{exhibitbox}[单票卡：源杰科技]
\small
\begin{tabularx}{\exhibitboxwidth}{>{\bfseries\raggedright\arraybackslash}p{2.5cm} >{\raggedright\arraybackslash\sloppy\hspace{0pt}}X}
\toprule
产业位置 & 激光芯片与高速光源 \\
当前价与市值 & CNY1776.91；市值约 CNY1526 亿元 \\
估值结论 & 综合目标 CNY1421.53，区间 CNY1109--1635，空间 -20.0\%，动作 \stance{riskamber}{中性观察} \\
估值锚点 & 内在锚 CNY434.08；市场锚 CNY1421.53；权重 内在47\% / 市场38\% / 券商15\% \\
核心催化 & 100G/200G EML 和硅光光源进入头部模块客户。 \\
失效条件 & 客户认证延迟、良率不稳，或双供应削弱议价能力。 \\
\bottomrule
\end{tabularx}
\end{exhibitbox}

投资含义上，上游战略属性最强之一，但当前价格已经前置长期兑现路径。 市场锚的作用是承认当前成交和共识定价已经给出的情绪溢价，但它不能替代利润、毛利率和现金流。组合处理应当把它放在与当前动作匹配的仓位层级，而不是因为光通信行业景气就无差别买入。下一季度最关键的是利润、毛利率和现金流是否同步，若只有收入增长而现金流或毛利率恶化，则说明订单质量低于股价隐含预期。


\subsection{亨通光电（600487）：光纤光缆、海缆与电力通信一体化}
亨通光电 位于本报告定义的 \textbf{光纤光缆} 层级，研究权重为 4\%。2026Q1 公司收入为 CNY177.91 亿元，归母净利润为 CNY11.05 亿元，毛利率 16.0\%，归母净利同比 +98.5\%。按照本报告季节性假设，2026E 归母净利为 CNY46.06 亿元，2027E EPS 为 CNY2.13。

\begin{exhibitbox}[单票卡：亨通光电]
\small
\begin{tabularx}{\exhibitboxwidth}{>{\bfseries\raggedright\arraybackslash}p{2.5cm} >{\raggedright\arraybackslash\sloppy\hspace{0pt}}X}
\toprule
产业位置 & 光纤光缆、海缆与电力通信一体化 \\
当前价与市值 & CNY110.81；市值约 CNY2684 亿元 \\
估值结论 & 综合目标 CNY86.43，区间 CNY67--99，空间 -22.0\%，动作 \stance{riskamber}{中性观察} \\
估值锚点 & 内在锚 CNY36.95；市场锚 CNY86.43；权重 内在38\% / 市场47\% / 券商15\% \\
核心催化 & 运营商/云侧光纤需求修复、海缆订单、营运资本改善。 \\
失效条件 & 电信资本开支走弱、光缆价格承压，或现金转化落后收入。 \\
\bottomrule
\end{tabularx}
\end{exhibitbox}

投资含义上，光纤光缆龙头兼具海缆和电网敞口，但 AI 数据中心纯度低于模块龙头。 市场锚的作用是承认当前成交和共识定价已经给出的情绪溢价，但它不能替代利润、毛利率和现金流。组合处理应当把它放在与当前动作匹配的仓位层级，而不是因为光通信行业景气就无差别买入。下一季度最关键的是利润、毛利率和现金流是否同步，若只有收入增长而现金流或毛利率恶化，则说明订单质量低于股价隐含预期。


\subsection{中天科技（600522）：光纤光缆、海缆与通信线缆}
中天科技 位于本报告定义的 \textbf{光纤光缆} 层级，研究权重为 4\%。2026Q1 公司收入为 CNY131.42 亿元，归母净利润为 CNY9.19 亿元，毛利率 15.6\%，归母净利同比 +46.4\%。按照本报告季节性假设，2026E 归母净利为 CNY38.29 亿元，2027E EPS 为 CNY1.27。

\begin{exhibitbox}[单票卡：中天科技]
\small
\begin{tabularx}{\exhibitboxwidth}{>{\bfseries\raggedright\arraybackslash}p{2.5cm} >{\raggedright\arraybackslash\sloppy\hspace{0pt}}X}
\toprule
产业位置 & 光纤光缆、海缆与通信线缆 \\
当前价与市值 & CNY61.29；市值约 CNY2071 亿元 \\
估值结论 & 综合目标 CNY47.81，区间 CNY37--55，空间 -22.0\%，动作 \stance{riskamber}{中性观察} \\
估值锚点 & 内在锚 CNY23.35；市场锚 CNY47.81；权重 内在38\% / 市场47\% / 券商15\% \\
核心催化 & 海缆交付、运营商光纤需求改善、毛利率修复。 \\
失效条件 & 电网/海缆项目延迟、光纤价格下降，或经营现金流恶化。 \\
\bottomrule
\end{tabularx}
\end{exhibitbox}

投资含义上，电力海洋与通信线缆平台提供盈利韧性，但不能按纯 AI 光模块估值。 市场锚的作用是承认当前成交和共识定价已经给出的情绪溢价，但它不能替代利润、毛利率和现金流。组合处理应当把它放在与当前动作匹配的仓位层级，而不是因为光通信行业景气就无差别买入。下一季度最关键的是利润、毛利率和现金流是否同步，若只有收入增长而现金流或毛利率恶化，则说明订单质量低于股价隐含预期。


\subsection{长飞光纤（601869）：光纤预制棒、光纤与光缆}
长飞光纤 位于本报告定义的 \textbf{光纤/预制棒} 层级，研究权重为 4\%。2026Q1 公司收入为 CNY36.95 亿元，归母净利润为 CNY4.95 亿元，毛利率 41.5\%，归母净利同比 +226.4\%。按照本报告季节性假设，2026E 归母净利为 CNY21.53 亿元，2027E EPS 为 CNY3.03。

\begin{exhibitbox}[单票卡：长飞光纤]
\small
\begin{tabularx}{\exhibitboxwidth}{>{\bfseries\raggedright\arraybackslash}p{2.5cm} >{\raggedright\arraybackslash\sloppy\hspace{0pt}}X}
\toprule
产业位置 & 光纤预制棒、光纤与光缆 \\
当前价与市值 & CNY543.00；市值约 CNY4481 亿元 \\
估值结论 & 综合目标 CNY434.40，区间 CNY339--500，空间 -20.0\%，动作 \stance{riskamber}{中性观察} \\
估值锚点 & 内在锚 CNY51.93；市场锚 CNY434.40；权重 内在38\% / 市场47\% / 券商15\% \\
核心催化 & 预制棒利用率修复、云/园区光纤建设、海外订单。 \\
失效条件 & 光纤供给过剩、ASP 下滑，或 Q2/Q3 利润未延续高 Q1 利润率。 \\
\bottomrule
\end{tabularx}
\end{exhibitbox}

投资含义上，预制棒/光纤/光缆一体化稀缺性较强，但数据中心需求需体现在利用率上。 市场锚的作用是承认当前成交和共识定价已经给出的情绪溢价，但它不能替代利润、毛利率和现金流。组合处理应当把它放在与当前动作匹配的仓位层级，而不是因为光通信行业景气就无差别买入。下一季度最关键的是利润、毛利率和现金流是否同步，若只有收入增长而现金流或毛利率恶化，则说明订单质量低于股价隐含预期。


\subsection{烽火通信（600498）：电信光网络设备与系统}
烽火通信 位于本报告定义的 \textbf{网络设备} 层级，研究权重为 3\%。2026Q1 公司收入为 CNY45.10 亿元，归母净利润为 CNY0.38 亿元，毛利率 26.4\%，归母净利同比 -30.4\%。按照本报告季节性假设，2026E 归母净利为 CNY1.75 亿元，2027E EPS 为 CNY0.16。

\begin{exhibitbox}[单票卡：烽火通信]
\small
\begin{tabularx}{\exhibitboxwidth}{>{\bfseries\raggedright\arraybackslash}p{2.5cm} >{\raggedright\arraybackslash\sloppy\hspace{0pt}}X}
\toprule
产业位置 & 电信光网络设备与系统 \\
当前价与市值 & CNY76.22；市值约 CNY975 亿元 \\
估值结论 & 综合目标 CNY53.35，区间 CNY42--61，空间 -30.0\%，动作 \stance{riskamber}{中性观察} \\
估值锚点 & 内在锚 CNY12.45；市场锚 CNY53.35；权重 内在45\% / 市场55\% / 券商0\% \\
核心催化 & OTN/ROADM 需求、运营商骨干网升级、盈利正常化。 \\
失效条件 & 运营商资本开支延迟、低毛利项目占比提高，或净利率未修复。 \\
\bottomrule
\end{tabularx}
\end{exhibitbox}

投资含义上，运营商光网络平台具备战略意义，但 Q1 利润分母薄，周期慢于云侧 AI capex。 市场锚的作用是承认当前成交和共识定价已经给出的情绪溢价，但它不能替代利润、毛利率和现金流。组合处理应当把它放在与当前动作匹配的仓位层级，而不是因为光通信行业景气就无差别买入。下一季度最关键的是利润、毛利率和现金流是否同步，若只有收入增长而现金流或毛利率恶化，则说明订单质量低于股价隐含预期。


\subsection{中兴通讯（000063）：网络设备、运营商/云基础设施与终端应用}
中兴通讯 位于本报告定义的 \textbf{网络设备/应用} 层级，研究权重为 6\%。2026Q1 公司收入为 CNY349.88 亿元，归母净利润为 CNY13.10 亿元，毛利率 28.3\%，归母净利同比 -46.6\%。按照本报告季节性假设，2026E 归母净利为 CNY59.57 亿元，2027E EPS 为 CNY1.37。

\begin{exhibitbox}[单票卡：中兴通讯]
\small
\begin{tabularx}{\exhibitboxwidth}{>{\bfseries\raggedright\arraybackslash}p{2.5cm} >{\raggedright\arraybackslash\sloppy\hspace{0pt}}X}
\toprule
产业位置 & 网络设备、运营商/云基础设施与终端应用 \\
当前价与市值 & CNY35.67；市值约 CNY1731 亿元 \\
估值结论 & 综合目标 CNY31.06，区间 CNY24--39，空间 -12.9\%，动作 \stance{riskamber}{中性} \\
估值锚点 & 内在锚 CNY31.03；市场锚 CNY25.68；权重 内在50\% / 市场35\% / 券商15\% \\
核心催化 & 运营商网络升级、AI 服务器/网络产品进展、利润率改善。 \\
失效条件 & 运营商资本开支走弱、海外政策压力，或非光通信业务稀释盈利质量。 \\
\bottomrule
\end{tabularx}
\end{exhibitbox}

投资含义上，覆盖运营商、云网络和应用层，但光通信只是业务组合的一部分。 市场锚的作用是承认当前成交和共识定价已经给出的情绪溢价，但它不能替代利润、毛利率和现金流。组合处理应当把它放在与当前动作匹配的仓位层级，而不是因为光通信行业景气就无差别买入。下一季度最关键的是利润、毛利率和现金流是否同步，若只有收入增长而现金流或毛利率恶化，则说明订单质量低于股价隐含预期。


\subsection{仕佳光子（688313）：PLC 分路器、AWG 与无源光芯片/器件}
仕佳光子 位于本报告定义的 \textbf{无源光芯片} 层级，研究权重为 3\%。2026Q1 公司收入为 CNY5.77 亿元，归母净利润为 CNY1.16 亿元，毛利率 34.1\%，归母净利同比 +24.7\%。按照本报告季节性假设，2026E 归母净利为 CNY5.05 亿元，2027E EPS 为 CNY1.40。

\begin{exhibitbox}[单票卡：仕佳光子]
\small
\begin{tabularx}{\exhibitboxwidth}{>{\bfseries\raggedright\arraybackslash}p{2.5cm} >{\raggedright\arraybackslash\sloppy\hspace{0pt}}X}
\toprule
产业位置 & PLC 分路器、AWG 与无源光芯片/器件 \\
当前价与市值 & CNY179.70；市值约 CNY812 亿元 \\
估值结论 & 综合目标 CNY125.79，区间 CNY98--145，空间 -30.0\%，动作 \stance{riskamber}{中性观察} \\
估值锚点 & 内在锚 CNY47.49；市场锚 CNY125.79；权重 内在47\% / 市场38\% / 券商15\% \\
核心催化 & AWG/PLC 器件在高速模块和电信/数据中心 WDM 系统中放量。 \\
失效条件 & 光芯片降价、客户认证延迟，或产能利用率低于 Q1 水平。 \\
\bottomrule
\end{tabularx}
\end{exhibitbox}

投资含义上，无源光芯片与 WDM/AWG/硅光生态相关性强，但估值需要高利用率支撑。 市场锚的作用是承认当前成交和共识定价已经给出的情绪溢价，但它不能替代利润、毛利率和现金流。组合处理应当把它放在与当前动作匹配的仓位层级，而不是因为光通信行业景气就无差别买入。下一季度最关键的是利润、毛利率和现金流是否同步，若只有收入增长而现金流或毛利率恶化，则说明订单质量低于股价隐含预期。


\subsection{长光华芯（688048）：高功率与数通激光芯片}
长光华芯 位于本报告定义的 \textbf{激光芯片} 层级，研究权重为 2\%。2026Q1 公司收入为 CNY1.30 亿元，归母净利润为 CNY0.04 亿元，毛利率 30.1\%，归母净利同比 +159.7\%。按照本报告季节性假设，2026E 归母净利为 CNY0.22 亿元，2027E EPS 为 CNY0.18。

\begin{exhibitbox}[单票卡：长光华芯]
\small
\begin{tabularx}{\exhibitboxwidth}{>{\bfseries\raggedright\arraybackslash}p{2.5cm} >{\raggedright\arraybackslash\sloppy\hspace{0pt}}X}
\toprule
产业位置 & 高功率与数通激光芯片 \\
当前价与市值 & CNY396.00；市值约 CNY698 亿元 \\
估值结论 & 综合目标 CNY277.20，区间 CNY216--319，空间 -30.0\%，动作 \stance{riskamber}{中性观察} \\
估值锚点 & 内在锚 CNY26.74；市场锚 CNY277.20；权重 内在47\% / 市场38\% / 券商15\% \\
核心催化 & 数通激光芯片认证、良率改善、高功率激光盈利修复。 \\
失效条件 & 研发转化延迟、良率不稳，或 Q2/Q3 盈利不能验证高倍数。 \\
\bottomrule
\end{tabularx}
\end{exhibitbox}

投资含义上，激光芯片期权高，但 Q1 利润分母仍小，基准情景要求快速认证和放量。 市场锚的作用是承认当前成交和共识定价已经给出的情绪溢价，但它不能替代利润、毛利率和现金流。组合处理应当把它放在与当前动作匹配的仓位层级，而不是因为光通信行业景气就无差别买入。下一季度最关键的是利润、毛利率和现金流是否同步，若只有收入增长而现金流或毛利率恶化，则说明订单质量低于股价隐含预期。


\subsection{光库科技（300620）：光纤器件、隔离器与铌酸锂/调制器期权}
光库科技 位于本报告定义的 \textbf{光器件} 层级，研究权重为 3\%。2026Q1 公司收入为 CNY4.26 亿元，归母净利润为 CNY0.45 亿元，毛利率 36.6\%，归母净利同比 +312.5\%。按照本报告季节性假设，2026E 归母净利为 CNY1.95 亿元，2027E EPS 为 CNY0.98。

\begin{exhibitbox}[单票卡：光库科技]
\small
\begin{tabularx}{\exhibitboxwidth}{>{\bfseries\raggedright\arraybackslash}p{2.5cm} >{\raggedright\arraybackslash\sloppy\hspace{0pt}}X}
\toprule
产业位置 & 光纤器件、隔离器与铌酸锂/调制器期权 \\
当前价与市值 & CNY362.43；市值约 CNY903 亿元 \\
估值结论 & 综合目标 CNY116.95，区间 CNY91--134，空间 -67.7\%，动作 \stance{riskred}{减持} \\
估值锚点 & 内在锚 CNY28.59；市场锚 CNY217.46；权重 内在50\% / 市场35\% / 券商15\% \\
核心催化 & 高速光器件附加值提升、薄膜铌酸锂商业化、海外需求修复。 \\
失效条件 & 调制器商业化延迟、毛利率压缩，或 AI 相关器件需求低于预期。 \\
\bottomrule
\end{tabularx}
\end{exhibitbox}

投资含义上，光器件和薄膜铌酸锂期权提高战略相关性，但股价已经计入较长兑现路径。 市场锚的作用是承认当前成交和共识定价已经给出的情绪溢价，但它不能替代利润、毛利率和现金流。组合处理应当把它放在与当前动作匹配的仓位层级，而不是因为光通信行业景气就无差别买入。下一季度最关键的是利润、毛利率和现金流是否同步，若只有收入增长而现金流或毛利率恶化，则说明订单质量低于股价隐含预期。


\subsection{通鼎互联（002491）：光纤光缆、通信线缆与网络集成}
通鼎互联 位于本报告定义的 \textbf{光纤光缆} 层级，研究权重为 3\%。2026Q1 公司收入为 CNY10.64 亿元，归母净利润为 CNY0.51 亿元，毛利率 27.0\%，归母净利同比 -61.1\%。按照本报告季节性假设，2026E 归母净利为 CNY2.23 亿元，2027E EPS 为 CNY0.20。

\begin{exhibitbox}[单票卡：通鼎互联]
\small
\begin{tabularx}{\exhibitboxwidth}{>{\bfseries\raggedright\arraybackslash}p{2.5cm} >{\raggedright\arraybackslash\sloppy\hspace{0pt}}X}
\toprule
产业位置 & 光纤光缆、通信线缆与网络集成 \\
当前价与市值 & CNY36.15；市值约 CNY444 亿元 \\
估值结论 & 综合目标 CNY26.03，区间 CNY20--30，空间 -28.0\%，动作 \stance{riskamber}{中性观察} \\
估值锚点 & 内在锚 CNY13.96；市场锚 CNY26.03；权重 内在45\% / 市场55\% / 券商0\% \\
核心催化 & 电信光纤/线缆需求恢复、通信线缆订单改善、现金流修复。 \\
失效条件 & 运营商资本开支延迟、线缆 ASP 承压，或 Q2/Q3 利润不能确认 Q1 修复。 \\
\bottomrule
\end{tabularx}
\end{exhibitbox}

投资含义上，属于完整产业链底座，但盈利质量和 AI 数据中心纯度低于模块龙头。 市场锚的作用是承认当前成交和共识定价已经给出的情绪溢价，但它不能替代利润、毛利率和现金流。组合处理应当把它放在与当前动作匹配的仓位层级，而不是因为光通信行业景气就无差别买入。下一季度最关键的是利润、毛利率和现金流是否同步，若只有收入增长而现金流或毛利率恶化，则说明订单质量低于股价隐含预期。


\subsection{永鼎股份（600105）：光缆、通信线缆与电信基础设施集成}
永鼎股份 位于本报告定义的 \textbf{光纤光缆} 层级，研究权重为 3\%。2026Q1 公司收入为 CNY12.46 亿元，归母净利润为 CNY1.59 亿元，毛利率 26.2\%，归母净利同比 -45.2\%。按照本报告季节性假设，2026E 归母净利为 CNY6.62 亿元，2027E EPS 为 CNY0.50。

\begin{exhibitbox}[单票卡：永鼎股份]
\small
\begin{tabularx}{\exhibitboxwidth}{>{\bfseries\raggedright\arraybackslash}p{2.5cm} >{\raggedright\arraybackslash\sloppy\hspace{0pt}}X}
\toprule
产业位置 & 光缆、通信线缆与电信基础设施集成 \\
当前价与市值 & CNY65.71；市值约 CNY961 亿元 \\
估值结论 & 综合目标 CNY46.00，区间 CNY36--53，空间 -30.0\%，动作 \stance{riskamber}{中性观察} \\
估值锚点 & 内在锚 CNY14.77；市场锚 CNY46.00；权重 内在45\% / 市场55\% / 券商0\% \\
核心催化 & 运营商线缆订单、电信基础设施交付、毛利率稳定。 \\
失效条件 & 项目延迟、线缆价格下行，或营运资本压力抵消利润修复。 \\
\bottomrule
\end{tabularx}
\end{exhibitbox}

投资含义上，线缆和电信基础设施修复可见，但业务慢周期、纯度低于 AI 光模块。 市场锚的作用是承认当前成交和共识定价已经给出的情绪溢价，但它不能替代利润、毛利率和现金流。组合处理应当把它放在与当前动作匹配的仓位层级，而不是因为光通信行业景气就无差别买入。下一季度最关键的是利润、毛利率和现金流是否同步，若只有收入增长而现金流或毛利率恶化，则说明订单质量低于股价隐含预期。


\subsection{长芯博创（300548）：光模块、有源光器件与相干传输组件}
长芯博创 位于本报告定义的 \textbf{光模块/光器件} 层级，研究权重为 4\%。2026Q1 公司收入为 CNY6.71 亿元，归母净利润为 CNY1.30 亿元，毛利率 49.2\%，归母净利同比 +45.0\%。按照本报告季节性假设，2026E 归母净利为 CNY5.42 亿元，2027E EPS 为 CNY2.34。

\begin{exhibitbox}[单票卡：长芯博创]
\small
\begin{tabularx}{\exhibitboxwidth}{>{\bfseries\raggedright\arraybackslash}p{2.5cm} >{\raggedright\arraybackslash\sloppy\hspace{0pt}}X}
\toprule
产业位置 & 光模块、有源光器件与相干传输组件 \\
当前价与市值 & CNY262.65；市值约 CNY759 亿元 \\
估值结论 & 综合目标 CNY189.11，区间 CNY148--217，空间 -28.0\%，动作 \stance{riskamber}{中性观察} \\
估值锚点 & 内在锚 CNY80.23；市场锚 CNY189.11；权重 内在55\% / 市场45\% / 券商0\% \\
核心催化 & 数通模块订单、相干产品增长、毛利率维持。 \\
失效条件 & 高速产品爬坡延迟、模块 ASP 承压，或毛利率回落。 \\
\bottomrule
\end{tabularx}
\end{exhibitbox}

投资含义上，光模块/器件平台具备数通敞口，但当前估值要求高速产品持续交付。 市场锚的作用是承认当前成交和共识定价已经给出的情绪溢价，但它不能替代利润、毛利率和现金流。组合处理应当把它放在与当前动作匹配的仓位层级，而不是因为光通信行业景气就无差别买入。下一季度最关键的是利润、毛利率和现金流是否同步，若只有收入增长而现金流或毛利率恶化，则说明订单质量低于股价隐含预期。


\subsection{德科立（688205）：相干光模块与传输设备组件}
德科立 位于本报告定义的 \textbf{相干/传输} 层级，研究权重为 3\%。2026Q1 公司收入为 CNY2.54 亿元，归母净利润为 CNY0.20 亿元，毛利率 25.7\%，归母净利同比 +35.1\%。按照本报告季节性假设，2026E 归母净利为 CNY0.89 亿元，2027E EPS 为 CNY0.68。

\begin{exhibitbox}[单票卡：德科立]
\small
\begin{tabularx}{\exhibitboxwidth}{>{\bfseries\raggedright\arraybackslash}p{2.5cm} >{\raggedright\arraybackslash\sloppy\hspace{0pt}}X}
\toprule
产业位置 & 相干光模块与传输设备组件 \\
当前价与市值 & CNY210.52；市值约 CNY345 亿元 \\
估值结论 & 综合目标 CNY73.50，区间 CNY57--85，空间 -65.1\%，动作 \stance{riskred}{减持} \\
估值锚点 & 内在锚 CNY40.10；市场锚 CNY134.73；权重 内在65\% / 市场35\% / 券商0\% \\
核心催化 & DCI/相干需求、海外客户拓展、盈利正常化。 \\
失效条件 & 相干订单延迟、客户集中度压力，或净利率不能放大。 \\
\bottomrule
\end{tabularx}
\end{exhibitbox}

投资含义上，补足 DCI 和运营商相干传输链条，但 Q1 利润基数仍偏小。 市场锚的作用是承认当前成交和共识定价已经给出的情绪溢价，但它不能替代利润、毛利率和现金流。组合处理应当把它放在与当前动作匹配的仓位层级，而不是因为光通信行业景气就无差别买入。下一季度最关键的是利润、毛利率和现金流是否同步，若只有收入增长而现金流或毛利率恶化，则说明订单质量低于股价隐含预期。


\subsection{腾景科技（688195）：光通信和激光用精密光学元件}
腾景科技 位于本报告定义的 \textbf{精密光学} 层级，研究权重为 3\%。2026Q1 公司收入为 CNY1.71 亿元，归母净利润为 CNY0.14 亿元，毛利率 37.2\%，归母净利同比 +10.7\%。按照本报告季节性假设，2026E 归母净利为 CNY0.63 亿元，2027E EPS 为 CNY0.57。

\begin{exhibitbox}[单票卡：腾景科技]
\small
\begin{tabularx}{\exhibitboxwidth}{>{\bfseries\raggedright\arraybackslash}p{2.5cm} >{\raggedright\arraybackslash\sloppy\hspace{0pt}}X}
\toprule
产业位置 & 光通信和激光用精密光学元件 \\
当前价与市值 & CNY187.60；市值约 CNY246 亿元 \\
估值结论 & 综合目标 CNY83.18，区间 CNY65--96，空间 -55.7\%，动作 \stance{riskred}{减持} \\
估值锚点 & 内在锚 CNY29.18；市场锚 CNY112.56；权重 内在55\% / 市场30\% / 券商15\% \\
核心催化 & 高端光学元件订单、激光/通信需求、毛利率改善。 \\
失效条件 & 订单波动、收入规模不足，或光学元件价格承压。 \\
\bottomrule
\end{tabularx}
\end{exhibitbox}

投资含义上，精密光学与高速光系统相关，但规模小，估值需要持续订单转化。 市场锚的作用是承认当前成交和共识定价已经给出的情绪溢价，但它不能替代利润、毛利率和现金流。组合处理应当把它放在与当前动作匹配的仓位层级，而不是因为光通信行业景气就无差别买入。下一季度最关键的是利润、毛利率和现金流是否同步，若只有收入增长而现金流或毛利率恶化，则说明订单质量低于股价隐含预期。


\subsection{联特科技（301205）：高速光模块与数通收发器}
联特科技 位于本报告定义的 \textbf{高弹性光模块} 层级，研究权重为 2\%。2026Q1 公司收入为 CNY2.13 亿元，归母净利润为 CNY0.03 亿元，毛利率 43.4\%，归母净利同比 -83.7\%。按照本报告季节性假设，2026E 归母净利为 CNY0.15 亿元，2027E EPS 为 CNY0.16。

\begin{exhibitbox}[单票卡：联特科技]
\small
\begin{tabularx}{\exhibitboxwidth}{>{\bfseries\raggedright\arraybackslash}p{2.5cm} >{\raggedright\arraybackslash\sloppy\hspace{0pt}}X}
\toprule
产业位置 & 高速光模块与数通收发器 \\
当前价与市值 & CNY315.86；市值约 CNY409 亿元 \\
估值结论 & 综合目标 CNY116.17，区间 CNY91--134，空间 -63.2\%，动作 \stance{riskred}{减持} \\
估值锚点 & 内在锚 CNY33.68；市场锚 CNY189.52；权重 内在55\% / 市场30\% / 券商15\% \\
核心催化 & 高速数通订单转化、客户认证、毛利率扩张。 \\
失效条件 & 认证延迟、Q2/Q3 利润不放量，或现金流不跟随收入。 \\
\bottomrule
\end{tabularx}
\end{exhibitbox}

投资含义上，高速模块弹性大，但 Q1 利润分母很小，目标价高度依赖执行。 市场锚的作用是承认当前成交和共识定价已经给出的情绪溢价，但它不能替代利润、毛利率和现金流。组合处理应当把它放在与当前动作匹配的仓位层级，而不是因为光通信行业景气就无差别买入。下一季度最关键的是利润、毛利率和现金流是否同步，若只有收入增长而现金流或毛利率恶化，则说明订单质量低于股价隐含预期。


\subsection{兆龙互连（300913）：高速数据线缆、通信线缆与铜互连}
兆龙互连 位于本报告定义的 \textbf{高速线缆/互连} 层级，研究权重为 2\%。2026Q1 公司收入为 CNY5.08 亿元，归母净利润为 CNY0.49 亿元，毛利率 21.1\%，归母净利同比 +49.0\%。按照本报告季节性假设，2026E 归母净利为 CNY2.02 亿元，2027E EPS 为 CNY0.70。

\begin{exhibitbox}[单票卡：兆龙互连]
\small
\begin{tabularx}{\exhibitboxwidth}{>{\bfseries\raggedright\arraybackslash}p{2.5cm} >{\raggedright\arraybackslash\sloppy\hspace{0pt}}X}
\toprule
产业位置 & 高速数据线缆、通信线缆与铜互连 \\
当前价与市值 & CNY50.75；市值约 CNY174 亿元 \\
估值结论 & 综合目标 CNY22.44，区间 CNY18--28，空间 -55.8\%，动作 \stance{riskred}{减持} \\
估值锚点 & 内在锚 CNY21.10；市场锚 CNY24.36；权重 内在59\% / 市场41\% / 券商0\% \\
核心催化 & 高速数据线缆需求、AI 服务器/网络互连订单、毛利率稳定。 \\
失效条件 & 铜互连需求不及预期、产品结构恶化，或价格竞争加剧。 \\
\bottomrule
\end{tabularx}
\end{exhibitbox}

投资含义上，高速线缆/铜互连补足机柜内和企业网络层，但不是纯光模块替代品。 市场锚的作用是承认当前成交和共识定价已经给出的情绪溢价，但它不能替代利润、毛利率和现金流。组合处理应当把它放在与当前动作匹配的仓位层级，而不是因为光通信行业景气就无差别买入。下一季度最关键的是利润、毛利率和现金流是否同步，若只有收入增长而现金流或毛利率恶化，则说明订单质量低于股价隐含预期。


\subsection{意华股份（002897）：高速连接器与通信互连组件}
意华股份 位于本报告定义的 \textbf{连接器/互连} 层级，研究权重为 2\%。2026Q1 公司收入为 CNY12.92 亿元，归母净利润为 CNY0.40 亿元，毛利率 20.0\%，归母净利同比 -37.7\%。按照本报告季节性假设，2026E 归母净利为 CNY1.83 亿元，2027E EPS 为 CNY1.10。

\begin{exhibitbox}[单票卡：意华股份]
\small
\begin{tabularx}{\exhibitboxwidth}{>{\bfseries\raggedright\arraybackslash}p{2.5cm} >{\raggedright\arraybackslash\sloppy\hspace{0pt}}X}
\toprule
产业位置 & 高速连接器与通信互连组件 \\
当前价与市值 & CNY84.50；市值约 CNY162 亿元 \\
估值结论 & 综合目标 CNY49.92，区间 CNY39--67，空间 -40.9\%，动作 \stance{riskred}{减持} \\
估值锚点 & 内在锚 CNY48.20；市场锚 CNY52.39；权重 内在59\% / 市场41\% / 券商0\% \\
核心催化 & 高速连接器订单、AI 服务器互连需求、盈利修复。 \\
失效条件 & 连接器需求放缓、非光通信业务稀释，或毛利率压缩。 \\
\bottomrule
\end{tabularx}
\end{exhibitbox}

投资含义上，连接器敞口属于数据中心互连链，但混合业务需要纯度折价。 市场锚的作用是承认当前成交和共识定价已经给出的情绪溢价，但它不能替代利润、毛利率和现金流。组合处理应当把它放在与当前动作匹配的仓位层级，而不是因为光通信行业景气就无差别买入。下一季度最关键的是利润、毛利率和现金流是否同步，若只有收入增长而现金流或毛利率恶化，则说明订单质量低于股价隐含预期。


\subsection{神宇股份（300563）：射频同轴线缆与高速通信互连}
神宇股份 位于本报告定义的 \textbf{线缆/互连} 层级，研究权重为 2\%。2026Q1 公司收入为 CNY2.04 亿元，归母净利润为 CNY0.13 亿元，毛利率 18.2\%，归母净利同比 +11.4\%。按照本报告季节性假设，2026E 归母净利为 CNY0.61 亿元，2027E EPS 为 CNY0.37。

\begin{exhibitbox}[单票卡：神宇股份]
\small
\begin{tabularx}{\exhibitboxwidth}{>{\bfseries\raggedright\arraybackslash}p{2.5cm} >{\raggedright\arraybackslash\sloppy\hspace{0pt}}X}
\toprule
产业位置 & 射频同轴线缆与高速通信互连 \\
当前价与市值 & CNY24.73；市值约 CNY48 亿元 \\
估值结论 & 综合目标 CNY12.95，区间 CNY10--18，空间 -47.7\%，动作 \stance{riskred}{减持} \\
估值锚点 & 内在锚 CNY13.70；市场锚 CNY11.87；权重 内在59\% / 市场41\% / 券商0\% \\
核心催化 & 高速通信线缆需求、客户结构改善、毛利率稳定。 \\
失效条件 & 低毛利产品占比提高、订单转化弱，或现金流落后利润。 \\
\bottomrule
\end{tabularx}
\end{exhibitbox}

投资含义上，线缆/互连拓宽产业链覆盖，但光通信纯度和 Q1 利润规模有限。 市场锚的作用是承认当前成交和共识定价已经给出的情绪溢价，但它不能替代利润、毛利率和现金流。组合处理应当把它放在与当前动作匹配的仓位层级，而不是因为光通信行业景气就无差别买入。下一季度最关键的是利润、毛利率和现金流是否同步，若只有收入增长而现金流或毛利率恶化，则说明订单质量低于股价隐含预期。
